% Generated draft: requires mathematical and source review.
\documentclass[11pt]{article}
\usepackage[margin=25mm]{geometry}
\usepackage{fontspec}
\setmainfont{Latin Modern Roman}
\usepackage{amsmath,amssymb,mathtools,unicode-math}
\setmathfont{Latin Modern Math}
\usepackage{enumitem,needspace,xurl,hyperref}
\setlength{\parindent}{0pt}
\setlength{\parskip}{5pt}
\setlength{\emergencystretch}{3em}
\newcommand{\N}{\mathbb{N}}
\newcommand{\Z}{\mathbb{Z}}
\newcommand{\Q}{\mathbb{Q}}
\newcommand{\R}{\mathbb{R}}
\newcommand{\C}{\mathbb{C}}
\newcommand{\F}{\mathbb{F}}
\newcommand{\A}{\mathbb{A}}
\newcommand{\PP}{\mathbb{P}}
\newcommand{\E}{\mathbb{E}}
\newcommand{\eps}{\varepsilon}
\newcommand{\dd}{\,\mathrm{d}}
\begin{document}
\section*{UnsolvedMath ID: 30006449}
\textbf{Source problem code:} catalogue entry
\subsection{Definitions and mathematical statement}
The supplied record is an incomplete theorem excerpt, marked \emph{unrecoverable} by the supplied statement verification. It does not determine a separately posed open problem. The following records its displayed identity and qualifications without promoting it to a universal conjecture.

The excerpt considers a suspended polynomial
\[
p(x,y,z,w)=x^2+y^2+g(z,w)
\]
and the threefold singularity \(X=(p=0,0)\). It denotes the corresponding Milnor fibre by \(F\) and the link by \(L\). In the usual interpretation for an isolated complex hypersurface germ, these are respectively a sufficiently small smooth local fibre and the intersection of the singular hypersurface with a sufficiently small sphere. This interpretation is conditional: the excerpt does not explicitly specify the ground field or the hypotheses on \(g\) ensuring that these objects have the intended properties.

A small \(\Q\)-factorialization is a proper birational morphism \(\pi:Y\to X\) whose exceptional locus has codimension at least two in \(Y\), with \(Y\) normal and \(\Q\)-factorial; the latter means that every Weil divisor has a positive integer multiple that is Cartier. For a small \(\Q\)-factorialization as considered in the excerpt, let \(F_1,\ldots,F_m\) denote the Milnor fibres of the resulting \(\Q\)-factorial singularities. Write \(SH^j(M)\) for degree-\(j\) symplectic cohomology of a Milnor fibre \(M\), and \(b_2(L)\) for the second Betti number of \(L\). The coefficient field, grading conventions, and finiteness assumptions needed for the dimensions are not specified in the supplied fragment.

The displayed assertion, reported there as Theorem 3, is
\[
\text{for every }k\in\Z\text{ with }k\leq 0,\qquad
\dim SH^{2k}(F)
=
\sum_{i=1}^{m}\dim SH^{2k}(F_i)+b_2(L).
\]
The excerpt presents this identity for suspended polynomials, but its opening refers to a refined version of an unidentified Conjecture 5.3 and to the singularities considered in an unidentified reference [1]. The supplied dated assessment restricts the reported proof to suspended polynomials in the considered quasihomogeneous compound-\(A\) families. Consequently, the complete hypotheses and scope of the displayed theorem cannot be recovered here, and the identity is not asserted for arbitrary suspended polynomials or all isolated compound Du Val singularities.

As additional context, the excerpt reports that the links of two quasihomogeneous compound-\(A_n\) singularities are contactomorphic if and only if the singularities are deformation equivalent in the sense of [1, Definition 5.1]. That definition and its precise scope are absent. The next theorem begins with Fermat-type equations
\[
p_i=x^2+y^2+z^{e_i}+w^{f_i}=0,
\]
but its conclusion and exponent hypotheses are missing. Neither this truncated theorem nor the contextual contactomorphism claim supplies a recoverable independent problem statement.
\subsection{Short English statement}
This record preserves a reported symplectic-cohomology identity relating a Milnor fibre to the Milnor fibres arising from a small \(\Q\)-factorialization and the second Betti number of its link. The supplied verification identifies the extraction as a proved-result fragment, with incomplete hypotheses, rather than a separately posed open problem. A precise universal conjecture cannot be recovered from the supplied evidence.
\subsection{Sources}
\begin{itemize}
\item Local UnsolvedMath registry: supplied complete record with permanent ID \(30006449\), exact source problem code \texttt{catalogue entry}, titled \emph{Proved Symplectic-Cohomology Factorization Fragment}, in the Oberwolfach Reports Open Problems collection. No proposer or individual author is supplied.
\item Supplied report citation: \emph{Singularities (2026)}, Oberwolfach Reports, pp.\ 2413--2490, DOI \texttt{10.4171/owr/2025/45}. The literature-source label is \emph{Oberwolfach Report 45/2025}. \url{https://doi.org/10.4171/owr/2025/45}.
\item Supplied literature source: \emph{Symplectic cohomology of quasihomogeneous cA singularities}. \url{https://doi.org/10.1007/s00209-026-03943-7}. The supplied record does not identify this citation explicitly with the excerpt's reference [1].
\item Literature assessment checked at \texttt{2026-08-22}: this is supplied dated curation, not independently verified. Its status label is \emph{partially solved}. It reports that the refined identity is proved for suspended polynomials in the considered quasihomogeneous compound-\(A\) families, while the extracted universal conjecture is not proved for all isolated compound Du Val singularities. The background explicitly requests independent verification of time-sensitive status and cited claims.
\item Supplied statement verification, not independently verified here: the record is marked \emph{unrecoverable}; the verification says that the official report presents the refined conjecture only to report its proof for the quasihomogeneous compound-\(A_n\) singularities considered there, and that it is not a separately posed open problem. No clean statement was promoted; the original record was retained for provenance.
\end{itemize}
\end{document}
